\chapter{Related Work}

In this chapter, we review frameworks for development of FPGA-based SOCs,
works that explore hardware-based optimizations for certain algorithms
and works that explore other ways of managing multiple extensions for RISC-V-based CPUs.

\section{RISC-V Extensions For Custom Instruction Integration}

In this section, we provide an overview of what the composable
custom extensions specification is and what it may be used for.

\subsection{CFU}

Composable Custom Extensions \cite{Gray2022CFU, Gray2025CXU, Gray2025CXUGithub} is a
RISC-V extension proposed by Jan Gray et al. It's an extension aimed at making custom
extensions modular and usable between different implementations of RISC-V CPUs.
The motivation behind the composable custom extensions is that many RISC-V CPUs
use some sort of custom extensions, yet there is no cross-implementation interoperability,
and so there is many implementations of a single extension (e. g. SIMD multiply-accumulate)
across multiple RISC-V CPU implementations. One solution to that is to provide a common logic
interface to connect extensions to, and that is exactly the Composable Custom Extensions RISC-V
ISA extension.

The first version of composable custom extensions if the Custom Function Unit (CFU).
It was proposed in 2020 by Jan Gray et al. and provides an interface to connect custom
extensions to along with timing, logic and combinatorial constraints for what an extension
may do \cite{Gray2022CFU}.
We will go over what constitutes a composable extension in the theoretical background chapter.

The CFU allows for a single extension to be used with a CPU,
and it requires the user to invoke a \textbf{custom-0|1|2|3} instruction,
even when not using assembly, such that when using a higher-level language,
the user would still need to invoke an assembly instruction.
The CFU also may not access memory. \cite{Gray2022CFU}

\subsection{CXU}

The CXU extension is a newer version of composable custom extensions.
It allows for up to $2^{\textbf{XLEN}}$ custom extensions.
It also allows for extensions to have a shared state and defines the limitations under which an extension may have access to memory. \cite{Gray2025CXU}

Notably, the CXU specification also defines the ABI, making it possible to call the extension from a higher-level language without invoking the assembly instruction directly.
An example of a CXU runtime is also provided, such that calls to the CXU extensions
would be wrapped in helper functions that select a proper extension, save the previous state
and clean up the context afterwards.
The specification was developed with support for operating systems in mind, and it
specifies how the OS and process context switching is to be done with respect to the CXU
state as well.

\subsection{Core-V-XIF}

\textbf{Core-V} E\textbf{X}tensible \textbf{I}nter\textbf{F}ace \cite{OpenHW2021XIF} is a
RISC-V extension interface for implementing custom ISA extensions and coprocessors.
As composable custom extensions, it's an interface that allows for custom extensions to be
integrated into any CPU that supports the Core-V-XIF interface.

Unlike composable custom extensions, Core-V-XIF allows for custom CSRs and overall is less
restrictive. Like the CFU Playgroud \cite{CFUPlayground2022}, there is an implementation of a
sandbox for custom extensions, Francis-V \cite{Egert2023FRANCISV}.
However, it's more tightly coupled to a certain CPU implementation, and would not allow
for straightforward interoperability of implemented extensions.

\section{Frameworks for FPGA-based SOC development}

In this section, we provide an overview for frameworks that provide a base for developing
FPGA-based SOCs. These frameworks usually provide a way to extend the SOC
with additional functionality, be it an ISA extension or an IO device, e. g.
a specialized USB port for debugging.
The frameworks usually consist of modules, which are then assembled into a bigger system.

\subsection{LiteX}

\textbf{LiteX} \cite{LiteXGitHub} is an open-source python-based framework for building
FPGA-based SoCs. It provides modules for the bus, IO and CPU integration
(notably VexiiRiscv \cite{VexiiRiscVGithub}).
It includes support for most popular dram modules used on FPGA devboards,
support for ethernet, PCIE and SATA buses, support for SD card readers,
inter-chip communication and most popular serial communication buses,
notably I2C, UART, SPI and USB.

The most popular implementation of an environment for development of CFUs,
the  CFU Playground \cite{CFUPlayground2022}, used LiteX as a basis for it's environment.

We use litex for our implementation as well.

\subsection{F4PGA}

\textbf{F4PGA} \cite{F4PGA2019} provides a fully open-source FPGA toolchain targeting multiple FPGA families. While not a SoC framework by itself,
it is commonly used together with LiteX or custom SoC
designs to enable reproducible, vendor-independent FPGA-based SoC development.

\subsection{MiSoC}

\textbf{MiSoC} \cite{MiSoC2014} is the predecessor of LiteX and introduced Python-based
SoC generation concepts.
Although deprecated, it is historically important and illustrates
early approaches to modular FPGA SoC construction, which directly influenced LiteX’s
architecture.

\subsection{PULP Platform}

\textbf{PULP} \cite{PULP2015} is an open-source RISC-V-based SoC platform focused
on energy-efficient parallel processing and accelerator integration.
It uses a cluster-based architecture and is widely used in academic
research for DSP, cryptography, and accelerator-centric SoC exploration.

\subsection{SpinalHDL}

\textbf{SpinalHDL} \cite{SpinalHDL2018} is a hardware construction language
often used to build extensible SoCs, most notably the VexRiscv core.
Its plugin-based CPU architecture enables fine-grained custom instruction
and pipeline extensions, making it highly relevant for composable extension research.

SpinalHDL is used in the \textbf{VexiiRISCV} \cite{VexiiRiscVGithub} implementation, and we use
VexiiRISCV as our CPU.

\section{Sandboxes for simplified development of custom instruction hardware}

\subsection{FRANCIS-V}

\textbf{FRANCIS-V}\cite{Egert2023FRANCISV}
(FRAmework for iNtegrating Custom InstructionS into RISC‑V systems) \cite{Egert2023FRANCISV} is a
framework aimed as making development and integration of custom extensions and coprocessors
easier while designing RISC-V processors.
It was proposed to make enginnering custom instructions easier.

The framework provices an interface and methodology for connecting custom instruction
hardware to a processor using the \textbf{Core-V-XiF} extension. It supports simulation,
compilation and verificcation workflows. It allows users to automate most parts of development,
much like the \textbf{CFU playground} and allows for customized RISC-V processors.

\begin{figure}[htb]
\centering
\begin{tikzpicture}[
    box/.style={rectangle, draw, rounded corners=3pt,
                minimum width=3.2cm, minimum height=0.75cm,
                text centered, font=\small},
    wbox/.style   = {rectangle, draw, dashed, rounded corners=4pt,
                     inner sep=8pt},
    arr/.style    = {->, >=stealth, thick},
    darr/.style   = {<->, >=stealth, thick, gray},
    node distance = 0.55cm and 1.3cm
]

\node[box, fill=purple!15]               (src)  {Custom Instr.\ RTL};
\node[box, fill=purple!10, below=of src] (gen)  {Code Generator};
\node[box, fill=purple!10, below=of gen] (drv)  {Driver / Wrapper};

\node[box, fill=teal!15,   right=2.8cm of src]  (xif)  {Core-V XIF Interface};
\node[box, fill=teal!10,   below=of xif]         (cpu)  {RISC-V CPU Core};
\node[box, fill=teal!10,   below=of cpu]         (sim)  {Simulation / Verilator};

\draw[arr] (src) -- (gen);
\draw[arr] (gen) -- (drv);

\draw[arr] (xif) -- (cpu);
\draw[arr] (cpu) -- (sim);

\draw[darr] (src)  -- (xif)  node[midway, above, font=\scriptsize]{};
\draw[darr] (drv)  -- (sim)  node[midway, below, font=\scriptsize]{};

\node[wbox, fit=(src)(gen)(drv), label={[font=\scriptsize]above:FRANCIS-V Tooling}] {};
\node[wbox, fit=(xif)(cpu)(sim), label={[font=\scriptsize]above:Target CPU Environment}] {};

\end{tikzpicture}
\caption{FRANCIS-V framework overview.}
\end{figure}

\subsection{CFU Playground}
The CFU Playground \cite{CFUPlayground2022} is an environment for development of custom
compute units on a litex SoC.

The CFU Playground \cite{CFUPlayground2022} is a full-stack open-source playground designed
for the rapid design and evaluation of Custom Function Units (CFUs)
to accelerate FPGA-based RISC-V soft CPUs.

Originally proposed for Tiny Machine Learning (tinyML) acceleration,
it utilizes the LiteX framework as its architectural base to provide a ready-made SoC environment
\cite{LiteXGitHub}. The playground allows users to implement specialized hardware logic and
write software scripts to invoke these extensions.

The architecture is based on the initial version of the Composable Custom Extension specification,
which connects a single extension to a hardware thread \cite{Gray2022CFU, CFUPlayground2022} and, usually
a single CPU.
While the interface supports up to 255 different function indices,
it is characterized by some constraints: the unit may not access memory directly,
and users must invoke instructions via assembly-level custom-0|1|2|3 opcodes,
even when programming in higher-level languages \cite{Gray2022CFU}.
These restrictions eventually led to the development of the more flexible CXU specification \cite{Gray2025CXU, Gray2025CXUGithub}.

\begin{figure}[htb]
\centering
\begin{tikzpicture}[
    box/.style={rectangle, draw, rounded corners=3pt,
                minimum width=3.2cm, minimum height=0.75cm,
                text centered, font=\small},
    wide/.style={rectangle, draw, rounded corners=3pt,
                minimum width=6.8cm, minimum height=0.75cm,
                text centered, font=\small},
    wbox/.style={rectangle, draw, dashed, rounded corners=4pt,
                inner sep=8pt},
    arr/.style={->, >=stealth, thick},
    darr/.style={<->, >=stealth, thick},
    node distance=0cm and 0cm
]
\node[wide, fill=blue!15] (user)  {User C / C++ Application};
\node[wide, fill=blue!10, below=of user] (cfu) {CFU Library (custom-0 intrinsics)};
\node[box, fill=green!15, below=of cfu, xshift=-2cm] (litex) {LiteX SoC (VexRiscv)};
\node[box, fill=green!10, right=of litex] (hw) {CFU Hardware (RTL)};
\node[wide, fill=orange!15, below=of litex, xshift=2cm] (fpga) {FPGA Bitstream};

\draw[decorate, decoration={brace, amplitude=6pt}, thick]
    ([xshift=-2.0cm]user.north west) -- ([xshift=-2.0cm]cfu.south west)
    node[midway, left=8pt, font=\scriptsize] {Software};
\draw[decorate, decoration={brace, amplitude=6pt}, thick]
    ([xshift=-2.0cm]litex.north west) -- ([xshift=-2.0cm]fpga.south west)
    node[midway, left=8pt, font=\scriptsize] {Hardware};

\end{tikzpicture}
\caption{CFU Playground full-stack architecture.}
\end{figure}

\section{Applications of Custom Instruction Extensions}

The interface standards and sandboxes above exist to serve a practical need: a wide
range of important algorithms benefit substantially from narrowly scoped hardware
acceleration. This section surveys the application domains most relevant to the case
studies in this thesis --- cryptography, data compression, and machine learning ---
and draws out the lessons that inform the design of the extensions evaluated in
Chapter~5.

\subsection{AES and Symmetric Cryptography}

AES \cite{NIST_AES} is based on a substitution-permutation network whose inner
operations --- SubBytes (an S-box lookup over $\text{GF}(2^8)$) and MixColumns
(matrix multiplication over $\text{GF}(2^8)$) --- map naturally onto
combinational hardware. Cui and Balasch \cite{Cui2023MaskingAES} demonstrate this
concretely in the context of masked AES, introducing constant-time finite-field
instructions that simultaneously improve performance and resistance to side-channel
attacks. Their result is especially relevant here because it shows that a small
number of carefully scoped ISA extensions --- rather than a monolithic AES
coprocessor --- can yield substantial gains while preserving software flexibility.

Primas et al. \cite{Primas2024NLUV} take a related approach with the NLU-V
extension, fusing S-box substitution and linear diffusion into a single hardware
block to minimise intermediate data movement for both AES and PRESENT. Both
works support the design choice made in Chapter~4 to expose $\text{GF}(2^8)$
arithmetic as a primitive CXU instruction rather than implementing a full AES round.

\subsection{Post-Quantum Cryptography}

Classical public-key cryptosystems are vulnerable to Shor's algorithm on
sufficiently powerful quantum computers, motivating the NIST post-quantum
standardisation effort \cite{Vranken2023UniformMul, Levi2020RISQV}. Lattice-based
schemes such as CRYSTALS-Kyber and NewHope rely on polynomial arithmetic,
particularly the Number Theoretic Transform (NTT), whose butterfly operations and
modular reductions are expensive on a standard ALU.

Sozbilir et al. \cite{Sozbilir2020Kyber} prototype NTT-specific ISA extensions on
VexRiscv and report up to an 85\% speedup for Kyber and NewHope, demonstrating
that the class of processor used in this work is a viable target for PQC
acceleration. Levi et al. \cite{Levi2020RISQV} extend this further with
RISQ-V, a set of 29 tightly coupled custom instructions achieving a 9.6$\times$
speedup for Kyber. Franchetti et al. \cite{Franchetti2024MultiWord} address the
related problem of multi-word carry-propagation arithmetic required by classical
public-key schemes, using a coprocessor interface structurally similar to the CXU.
Vranken et al. \cite{Vranken2023UniformMul} argue for a uniform set of
multiply-accumulate and modular-reduction primitives capable of serving both
classical and post-quantum schemes, which foreshadows the approach taken for the
MAC extension in Chapter~5.

These works collectively motivate extending the CXU Playground to post-quantum
workloads as a natural direction for future work.

\subsection{Hashing}

Alawadi et al. \cite{Alawadi2025SHA3} implement the full Keccak-f permutation as a
single custom RISC-V instruction, reporting substantial reductions in instruction
count and cycle time for SHA-3. Their area-performance trade-off analysis provides a
useful reference for understanding the cost of integrating high-complexity primitives
into the CXU.

\subsection{Data Compression}

GZIP/DEFLATE combines LZ77 back-reference encoding with Huffman entropy coding
\cite{RFC1951}. The performance bottleneck lies in the inner decompression loop,
which requires repeated bit-field extractions, variable-width mask generation, and
modular address arithmetic --- all operations that reduce to a handful of hardware
gates but require several ALU instructions in software. The custom instructions
designed in Section~3.9 and evaluated in Chapter~5 target exactly this hot path.
To our knowledge, no prior work has applied the CXU or CFU mechanism to GZIP
acceleration, making this a novel case study.

\subsection{Machine Learning Inference}

Neural network inference on microcontrollers is dominated by Multiply-Accumulate
(MAC) operations in convolutional and fully connected layers. TensorFlow Lite Micro
\cite{David2021TFLiteMicro} provides portable INT8-quantised inference kernels for
this class of device. The CMSIS-NN library \cite{Lai2018CMSISNN} shows that ARM DSP
SIMD MAC instructions yield up to a 4.6$\times$ improvement over scalar code; the
PULP platform \cite{PULP2015} achieves comparable gains on RISC-V with
hardware-loop and dot-product extensions. The MAC extension evaluated in
Chapter~5 follows the same principle, exposing a fused multiply-accumulate
operation as a single-cycle CXU instruction.

MiniFloat-NN and ExSdotp \cite{Oliveira2022MiniFloat} push this further with
float8 dot-product accumulation for low-precision training, demonstrating that the
same CXU-style integration pattern scales to more exotic arithmetic formats.

\subsection{Linux on LiteX}

Running benchmarks inside a full operating system is important for obtaining
performance numbers that reflect real-world deployment conditions rather than
isolated microbenchmarks. The Linux-on-LiteX-VexRiscv project \cite{LinuxOnLiteXGitHub}
establishes the feasibility of booting Linux on a soft RISC-V core in a LiteX SoC,
and documents the required hardware configuration (MMU, atomic instructions,
supervisor privilege mode). This work adopts the same configuration and extends it
with CXU hardware.
